\section{中心三次量子化からのambientの内在的復元}
\label{sec:pg-ambient-recovery}
本節の二次包絡の主張は三次超曲面に特有でなく、
任意次元・任意の関係次数$d\ge3$について成立する。
本文と同じく、周囲代数を$A$、その中心正則商を$B$と書く。
本節では、系\ref{cor:pg-equivariant-character-descent}の量子化族を含め、
中心商の同型類の判定に用いる一般原理を証明する。
この復元原理だけから、固定した非零変形パラメータでの
非可換代数の同型類の像の次元は結論しない。

\begin{theorem}[任意次元のquadratic envelope]\label{thm:pg-ambient-recovery-general}
$k$を体、$V$を有限次元$k$-ベクトル空間とし、
次数一生成の二次連結次数付き代数
\begin{equation}\label{eq:pg-ambient-B}
 A=T(V)/(R_2)
\end{equation}
を考える。$G\in A_d$を非零斉次元、$d\ge3$とし、
$G$の生成する両側イデアルが$(G)=AG=GA$であると仮定する。
$B=A/(G)$と置くと、
\begin{equation}\label{eq:pg-ambient-envelope}
 A_{\mathrm{can}}(B):=
 \frac{T(B_1)}{\left\langle
 \ker(B_1\otimes B_1\longrightarrow B_2)\right\rangle}
 \simeq A
\end{equation}
であり、また
\begin{equation}\label{eq:pg-ambient-line}
 \ker\bigl((A_{\mathrm{can}}(B))_d
 \longrightarrow B_d\bigr)=kG
\end{equation}
となる。この復元は次数付き代数同型に関して関手的である。
$G$の正則性、中心性、$A$のAS正則性は
二つの復元式自体には必要ない。
\end{theorem}
\begin{proof}
$(G)$は$d$次以上にしか成分を持たず、特に
\begin{equation}\label{eq:pg-ambient-low}
 B_1=A_1=V,\qquad B_2=A_2
\end{equation}
である。よって$B_1\otimes B_1\to B_2$の核は$R_2$に等しい。
$A$は二次関係のみで表示されるので最初の同型が従う。
$(G)_d=A_0G=GA_0=kG$より第二式も従う。
両構成は積と商写像のみによるから同型に関して自然である。
\end{proof}

\begin{theorem}[偏極付き圏と$\mathbb Z$-代数からの復元]
\label{thm:pg-ambient-recovery-polarized}
上の定理の仮定に加え、$B$は左右noetherianな
連結次数付きAS Gorenstein代数でinjective dimensionが
$2$以上であるとする。このとき
\begin{equation}\label{eq:pg-ambient-section}
 B\xrightarrow{\sim}
 \bigoplus_{n\ge0}
 \operatorname{Hom}_{\operatorname{qgr}B}
   (\mathcal O_B,\mathcal O_B(n))
\end{equation}
は積を保つ同型である。
従って偏極付き三つ組
\begin{equation}\label{eq:pg-ambient-triple}
 (\operatorname{qgr}B,\mathcal O_B,(1))
\end{equation}
から$A$と直線$kG$を復元できる。
具体的には$E_i=\mathcal O_B(-i)$として
\begin{equation}\label{eq:pg-ambient-Zalg}
 \mathcal Z_{ij}=
 \operatorname{Hom}_{\operatorname{qgr}B}(E_j,E_i)
 \qquad (i\le j)
\end{equation}
の次数一合成の二次核のみを課した
$\mathbb Z$-代数が$A$の付随$\mathbb Z$-代数となる。
さらに$A$が左右noetherianなら
$\operatorname{qgr}B\to\operatorname{qgr}A$は
完全忠実なAbel関手である。
\end{theorem}
\begin{proof}
AS Gorenstein性とinjective dimensionの仮定により
\begin{equation}\label{eq:pg-ambient-extvan}
 \operatorname{Ext}_B^i(k,B)=0
 \qquad (i=0,1)
\end{equation}
となる。有限次元次数付き加群$L$は$k$のシフトによる
有限filtrationを持つので、
$\operatorname{Hom}_{\operatorname{gr}B}(L,B(n))$と
$\operatorname{Ext}^1_{\operatorname{gr}B}(L,B(n))$は消滅する。
$0\to B_{\ge r}\to B\to B/B_{\ge r}\to0$にこれを適用し、
商圏Homの直極限表示を用いれば
式\eqref{eq:pg-ambient-section}を得る。
積の対応は次数シフトを介した合成の定義による。
このsection algebraの原理はArtin--Zhangの理論
\cite{PAZAmbient}の本仮定下での直接的帰結である。

従って$\mathcal Z_{ij}=B_{j-i}$であり、
各$i$での長さ二の合成の核は
$\ker(B_1\otimes B_1\to B_2)$と一致する。
定理\ref{thm:pg-ambient-recovery-general}により
二次包絡は$A$の付随$\mathbb Z$-代数となる。
最後に$B$-加群を$A$-加群とみなす制限関手は完全忠実で、
有限次元torsionの条件も一致するため商Abel圏に降りる。
この最後の完全忠実性は導来圏での完全忠実性を意味しない。
\end{proof}

\begin{corollary}[一般次元の中心三次量子化への適用]
\label{cor:pg-ambient-recovery-quantization}
第\ref{sec:pg-all-dimensions}節の量子化で得られる、
$N$生成元の二次Koszul AS正則ambient $A_\eta$と
中心正則三次元$\Omega_\eta$の商
$B_\eta=A_\eta/(\Omega_\eta)$について、
定理\ref{thm:pg-ambient-recovery-general}が成立する。
さらに各ファイバー$B_\eta$が上のnoetherian
AS Gorenstein仮定を満たす範囲では、
その偏極付き三つ組からambientと中心三次直線が復元される。
この結論は各ファイバーの内在的同型類についての主張であり、
全Poisson moduli stackからnc cubic moduli stackへの
大域的量子化射を構成したことを意味しない。
\end{corollary}
\begin{proof}
量子化定理のambientは二次で、
中心三次元は正則であるから最初の定理が適用できる。
第\ref{sec:pg-all-dimensions}節で示されたAS Gorenstein性と、
所定のnoetherian条件から第二の定理を適用する。
\end{proof}

\subsection{Serre関手による復元と一般次元での限界}
\label{sec:pg-ambient-serre}
構造対象だけを指定して次数シフトを忘れる復元は、
上記の一般定理より強い問題である。
例えば$N$生成元の三次中心商について
滑らかさとCalabi--Yau型Serre双対性の仮定の下では、
\begin{equation}\label{eq:pg-ambient-serre-r}
 \mathsf S\simeq(3-N)[N-2],\qquad
 \mathsf S^{-1}[N-2]\simeq(N-3)
\end{equation}
となる。この式だけで内在的に得られるのは
$(N-3)$-Veroneseであり、次数一シフトそのものとは限らない。
一般$N$についてVeroneseから元の代数を復元できるとは
ここでは主張しない。またNakayama自己同型が非自明なら
Serre関手にその捻りが加わるため、式
\eqref{eq:pg-ambient-serre-r}の使用には追加仮定が必要である。

\begin{proposition}[二変数中心による偶数Veronese復元]
\label{prop:pg-ambient-even}
$k=\mathbb C$とする。$B$を連結次数付き代数とし、
$S=k[u,v]\subset Z(B)$は次数一の中心多項式環、
$B$は$S$上自由、$Z(B)_2=S_2$と仮定する。
このとき偶数Veronese $R=B^{(2)}$の次数付き同型類から
$B$の次数付き同型類を復元できる。
$S$の平方根に関する二つの持上げは
奇数次数への符号作用で異なる。
\end{proposition}
\begin{proof}
$z\in Z(R)$とし、$a\in B$を奇数次数の元とする。
$ua$は偶数次数であるため、
\begin{equation}\label{eq:pg-ambient-even-center-comm}
 u[z,a]=[z,ua]=0
\end{equation}
である。$B$は$S$上自由なので$u$は正則であり、
$z$は全ての$a$と可換である。従って
\begin{equation}\label{eq:pg-ambient-even-center}
 Z(R)=Z(B)^{(2)},\qquad
 Z(R)_1=S_2
\end{equation}
を得る。$Z(R)_1$の生成する中心部分代数は
\begin{equation}\label{eq:pg-ambient-conic}
 D=k[u^2,uv,v^2]
 \simeq k[a,b,c]/(ac-b^2)
\end{equation}
である。具体的には$D_1$の基底$a,b,c$を選び、唯一の二次関係を
$ac-b^2=0$とする。その二次Veronese表示
$a=u^2$、$b=uv$、$c=v^2$により$S=k[u,v]$を取る。
この選択は復元した代数の同型類に影響しない。
実際、$D$の次数付き自己同型は滑らかなconic
$\operatorname{Proj}D\simeq\mathbb P^1$の自己同型を誘導し、
$\mathcal O_{\mathbb P^1}(2)$を保つ。
従って$k=\mathbb C$上では$\mathrm{GL}_2$の元へ持ち上がり、
持上げの核は共通符号$\{\pm1\}$である。
以下では$R,D$の次数を二倍し、$S$の元の次数と合わせる。
$D$の$S$への中心的拡大を使い、
\begin{equation}\label{eq:pg-ambient-even-localizations}
 T_u=R[(u^2)^{-1}]
       \otimes_{D[(u^2)^{-1}]}S[u^{-1}],
 \quad
 T_v=R[(v^2)^{-1}]
       \otimes_{D[(v^2)^{-1}]}S[v^{-1}]
\end{equation}
と置く。$S[u^{-1}]=D[(u^2)^{-1}]\oplus uD[(u^2)^{-1}]$であり、
$B[u^{-1}]$の奇数部分は$u$とその偶数部分の積に等しい。
従って偶奇の次数を分けた積写像により
$T_u\simeq B[u^{-1}]$、$T_v\simeq B[v^{-1}]$となる。
$B$は$S$上自由であり、$k[u,v]$の高さ二の原点を
除いた正則関数は多項式なので、
\begin{equation}\label{eq:pg-ambient-even-intersection}
 B[u^{-1}]\cap B[v^{-1}]=B
\end{equation}
が共通局所化内で成り立つ。従って$B$を復元する。
二つの平方根の差は$u,v\mapsto-u,-v$であり、
次数の偶奇に関する符号に一致する。
\end{proof}


\begin{lemma}[五生成元Jacobian族の次数二中心]
\label{lem:pg-ambient-center2-Jacobian}
$F\in\mathbb C[x,y,z,u,v]_3$とし、
$u,v$を中心変数として
\begin{equation}\label{eq:pg-ambient-Jacobian-rel}
 [x,y]=\operatorname{sym}(F_z),\quad
 [y,z]=\operatorname{sym}(F_x),\quad
 [z,x]=\operatorname{sym}(F_y)
\end{equation}
で定める二次代数を$A_F$とする。
$f_3=F(x,y,z,0,0)$は滑らかな平面三次を定め、
$x,y,z$-次数のPBW濾過について
\begin{equation}\label{eq:pg-ambient-Jacobian-gr}
 \operatorname{gr}A_F\simeq S_{f_3}\otimes\mathbb C[u,v],
 \qquad Z(S_{f_3})_1=Z(S_{f_3})_2=0
\end{equation}
と仮定する。さらに$G_F\in Z(A_F)_3$は
最高$x,y,z$-次数記号が非零な中心正則元とする。
$B_F=A_F/(G_F)$について
\begin{equation}\label{eq:pg-ambient-Jacobian-center2}
 Z(A_F)_2=Z(B_F)_2=\mathbb C[u,v]_2
\end{equation}
が成立する。
\end{lemma}
\begin{proof}
PBW濾過で中心元の最高記号を取れば、
$S_{f_3}$の一次、二次中心の消滅から
$Z(A_F)_2=\mathbb C[u,v]_2$を得る。
$B_F$の次数二中心元を$Q\in(A_F)_2$へ持ち上げる。
中心性と三次の商核が$\mathbb CG_F$であることから
\begin{equation}\label{eq:pg-ambient-Jacobian-Qcomm}
 [Q,x]=c_xG_F,\quad[Q,y]=c_yG_F,\quad
 [Q,z]=c_zG_F
\end{equation}
と書ける。$G_F$は中心的なのでJacobi恒等式は
$[Q,[x,y]]=[Q,[y,z]]=[Q,[z,x]]=0$を与える。
二次関係へ代入し、中心性と正則性によって
\begin{equation}\label{eq:pg-ambient-Jacobian-Hessian}
 \sum_{i=x,y,z}c_i\,\partial_i\partial_jF=0,
 \qquad j=x,y,z
\end{equation}
を得る。$u=v=0$で制限すれば
$\partial_j(\sum_i c_i\partial_i f_3)=0$となる。
括弧内は斉次二次式ゆえ零である。
非零の定数ベクトル$c$についてこの等式が成立すれば、
$f_3$は一次方向に依存しない錐であり、
滑らかな平面三次であることに反する。
従って$c_x=c_y=c_z=0$。
$Q$は$A_F$の中心元であるから
$Q\in\mathbb C[u,v]_2$となる。
\end{proof}

\begin{corollary}[三次三様体の構造対象付き導来復元]
\label{cor:pg-ambient-derived-threefold}
$k=\mathbb C$とし、$A$を左右noetherianな五次元Koszul
AS正則Calabi--Yau代数とする。
$G\in Z(A)_3$を正則中心元、$B=A/(G)$とする。
$\operatorname{qgr}B$は有限大域次元を持ち、そのSerre関手が
$\mathsf S\simeq(-2)[3]$であると仮定する。
さらに$S=k[u,v]\subset Z(B)$が次数一中心多項式環で、
$B$は$S$上自由、$Z(B)_2=S_2$であるとする。
すると
\begin{equation}\label{eq:pg-ambient-derived-input}
 (D^b(\operatorname{qgr}B),\mathcal O_B)
\end{equation}
だけから$B$の次数付き同型類、ambient $A$、
中心三次直線$kG$を復元できる。
次数シフト$(1)$は入力として指定しなくてよい。
\end{corollary}
\begin{proof}
Serre関手は内在的なので、
$\mathsf S^{-1}[3]\simeq(2)$から
\begin{equation}\label{eq:pg-ambient-derived-Veronese}
 \bigoplus_{n\ge0}
 \operatorname{Hom}(\mathcal O_B,\mathsf S^{-n}[3n]\mathcal O_B)
 \simeq B^{(2)}
\end{equation}
が得られる。ここでは、$B$が左右noetherianなAS Gorenstein代数で
injective dimension $4$を持つことから、
定理\ref{thm:pg-ambient-recovery-polarized}のsection algebra同型を用いた。
命題\ref{prop:pg-ambient-even}で$B$を復元し、
定理\ref{thm:pg-ambient-recovery-general}を適用する。
\end{proof}

\begin{remark}[既存の三次三様体族と未解決部分]
Jacobian型五生成元族については、
純三変数の滑らかな三次式の低次数中心の消滅と
Hessianの非退化性を用い、一般点で
$Z(B)_2=k[u,v]_2$を示す補題がある。
$B$の$k[u,v]$上の自由性には、ambientの相対PBW性に加えて
商の最高active記号を確認する必要がある。
例えば$\operatorname{gr}A=S_{f_3}\otimes k[u,v]$で、
$G$の最高active記号が$S_{f_3}$の正則な三次元$g_3$ならば、
\begin{equation}\label{eq:pg-ambient-quotient-symbol}
 \operatorname{gr}(A/(G))\simeq(S_{f_3}/(g_3))\otimes k[u,v]
\end{equation}
となる。右辺の各active次数の$k$-基底を持ち上げ、
最高記号について帰納することで$B$の$k[u,v]$-自由基底を得る。
ambientが自由であるという条件だけでは商の自由性は従わず、
例えば中心元$u^3$による商は$k[u,v]$上torsionを持つ。
ただしこの事実だけで$\operatorname{qgr}B$の
有限大域次元が全閉点で保証されるとはいわない。
これらの追加条件なしに一般次元でSerre復元を主張しない。
また構造対象を忘れた任意の導来同値、
Kuznetsov成分のみからのambient復元は結論しない。
\end{remark}
